\documentclass[11pt]{exam}
\usepackage[utf8]{inputenc}
\usepackage[a4paper,top=2cm,left=2cm,right=2cm,bottom=2cm]{geometry}
\usepackage{import}
\usepackage[T1]{fontenc}
\usepackage[german,ngerman]{babel}
\usepackage[table]{xcolor}
\usepackage{amsthm, esvect}
\usepackage{mathtools}
\RequirePackage{amssymb, amsfonts, amsmath, latexsym, verbatim, xspace}
\usepackage{systeme}
\usepackage{multicol}
\usepackage{multirow}
\usepackage{framed}
\usepackage{array}
\usepackage{ragged2e}
\usepackage{graphicx}
\usepackage{pdflscape}
\usepackage{epstopdf}
\usepackage{booktabs}
\usepackage{pdfpages}
\usepackage{float} % Allows putting an [H] in \begin{figure} to specify the exact location of the figure
\usepackage{wrapfig} % Allows in-line images such as the example fish picture
\usepackage{tikz, graphicx} % tikz
\usepackage[position=top]{subfig}
\usepackage[bookmarks=true]{hyperref}%Links
\usepackage{enumerate}
\usepackage{enumitem}
\usepackage[german,ruled,vlined,linesnumbered,algosection]{algorithm2e}
\usepackage[labelfont=sf]{caption}
\usepackage{lmodern}
\usepackage{lscape}
\usepackage{tabularx}
\usepackage{systeme}
\usepackage{hhline}
\usepackage{subfiles}
\usepackage{stackengine}
\usepackage{setspace}
\usepackage{MnSymbol}
\usepackage{tcolorbox}
\usepackage{bbm}
\usepackage{url}
\usepackage{xparse}
\usepackage{sectsty}
\usepackage{array}
\usepackage{ragged2e}
\usepackage{listings} %Stellt Code-Snippets sauber dar

% definiert die Sprache
\selectlanguage{German}

% add command for different dates in Swiss fashion
\newcommand{\monthwordlong}[1]{\ifcase#1 \or Januar\or Februar\or März\or April\or Mai\or Juni\or Juli\or August\or September\or Oktober\or November\or Dezember\fi}
\newcommand{\monthwordshort}[1]{\ifcase#1\or Jan\or Feb\or Mar\or Apr\or Mai\or Jun\or Jul\or Aug\or Sept\or Okt\or Nov\or Dez\fi}
\newcommand{\datelong}{\number\day.\:\monthwordlong{\the\month}\:\number\year}
\newcommand{\dateshort}{\number\day.\:\monthwordshort{\the\month}\:\number\year}

% define color of links inside a text
\definecolor{linkblue}{RGB}{0, 0, 80}
\definecolor{plotblue}{RGB}{100, 100, 255}
\definecolor{myblue}{RGB}{8,164,214}
\definecolor{myred}{RGB}{143,42,41}
\definecolor{forestgreen}{RGB}{34,139,34}
\definecolor{highdive}{RGB}{10,169,194}
\definecolor{charcoal}{RGB}{74,74,74}
\hypersetup{
	colorlinks=true,
	linkcolor=linkblue,
	filecolor=magenta,
	urlcolor=plotblue,
	citecolor=linkblue,
	pdfpagemode=FullScreen,
}


\lstdefinestyle{mystyle}{ % definiert den style des Code-Snippets
	%backgroundcolor=\color{backcolour},
	%commentstyle=\color{codegreen},
	%keywordstyle=\color{magenta},
	%numberstyle=\tiny\color{codegray},
	%stringstyle=\color{codepurple},
	basicstyle=\ttfamily\footnotesize,
	breakatwhitespace=false,
	breaklines=true,
	captionpos=b,
	keepspaces=true,
	%numbers=left,
	%numbersep=5pt,
	showspaces=false,
	showstringspaces=false,
	showtabs=false,
	tabsize=4
}
\lstset{style=mystyle}
\qformat{\Large\textbf{Aufgabe \thequestion\hfill\thepoints}}
\newlist{partes}{enumerate}{3}
\setlist[partes]{label=(\alph*)}
\newcommand{\parte}{\item}
\newlist{subpartes}{enumerate}{3}
\setlist[subpartes]{label=\roman*)}
\newcommand{\subparte}{\item}
\setlist[enumerate,1]{%(
	leftmargin=*, itemsep=12pt, label={\textbf{\arabic*.)}}}
\newcommand\answerbox{%%
	\fbox{\rule{2in}{0pt}\rule[-0.1ex]{0pt}{4ex}}}
\pointpoints{Punkt}{Punkte}
\hqword{Frage:}
\hpword{Mögliche Punkte:}
\hsword{Erreicht:}
\htword{Total}
\renewcommand*\half{.5}
\renewcommand{\solutiontitle}{\noindent\textbf{Lösung:}\par\noindent}

% Graphed boxes
\newcommand{\karos}[2]{
  \begin{tikzpicture}
    \draw[step=0.4cm,color=cyan,dotted] (0,0) grid (#1 cm ,#2 cm);
  %  \draw((#1)/2 cm,0)--((#1)/2 cm,/#2)/2cm)
  \end{tikzpicture}
}
\usetikzlibrary{arrows, petri, topaths, graphs}
\linespread{1.2} % Line spacing

%=========================================================
\newenvironment{parcols}[1][]{
    \begin{parts}\begin{multicols}{#1}
}{
    \end{multicols}\end{parts}
}
%=========================================================
\newenvironment{itemcols}[1][]{
	\begin{itemize}\begin{multicols}{#1}
		}{
	\end{multicols}\end{itemize}
}

%=========================================================
\newcommand{\antwort}[1]{
	{\setlength{\textwidth}{\textwidth}\fillwithgrid{#1cm}\vspace{5pt}}
}

%=========================================================
% Karriertes Papier für Antwort MIT OVERLAY
\tcbuselibrary{breakable,skins}
\usetikzlibrary{mindmap,trees}
\newlength\tindent
\setlength{\tindent}{\parindent}
\setlength{\parindent}{0pt}
\setlength{\parskip}{0pt}

\newtcolorbox{gridspace}[1]{%
	enhanced,
	breakable,
	blanker,
	boxrule=0pt,
	left=5mm,
	right=5mm,
	top=5.75mm,
	bottom=0mm,
	boxsep=0pt,
	height=#1cm,
	width=16cm,
	before upper={
		\begingroup
		\setlength{\baselineskip}{5.75mm}
		\setlength{\parskip}{0pt}
		\setlength{\lineskip}{0pt}
		% keep every paragraph on the grid rhythm
		%\everypar{\vspace{0pt}\strut}
		% math formulas shifted by one grid row
		\let\oldbracketopen\[
		\renewcommand{\[}{\vspace{-5mm}\oldbracketopen}
		% Align displayed formulas to the grid
		\setlength{\abovedisplayskip}{0pt}
		\setlength{\belowdisplayskip}{0pt}
		\setlength{\abovedisplayshortskip}{0pt}
		\setlength{\belowdisplayshortskip}{0pt}
		% Enumerate follows the same rhythm
		\setlist[enumerate]{leftmargin=6mm,itemsep=0mm,topsep=0pt,parsep=0pt,partopsep=0pt}
		\setlist[itemize]{leftmargin=6mm,itemsep=10mm,topsep=0pt,parsep=0pt,partopsep=0pt}},
	after upper={
		\endgroup},
	underlay={
		\draw[step=5mm,color=gray!55] (interior.south west) grid (interior.north east);}
}

\pagestyle{headandfoot}
\firstpageheader{}{}{\teacher}
\runningheader{}{}{\teacher}
\firstpagefooter{}{}{Seite \thepage\:von \numpages}
\runningfooter{}{}{Seite \thepage\:von \numpages}
%%
% Math Arrows
%%
\newcommand{\same}{\ensuremath{\Longleftrightarrow}}
\renewcommand{\to}{\ensuremath{\longrightarrow}}
\newcommand{\qsame}{\quad\same\quad}
\newcommand{\dsame}{\ensuremath{:\Leftrightarrow}}
\newcommand{\qdsame}{\quad\dsame\quad}
\newcommand{\ldsame}{\ensuremath{:\Longleftrightarrow}}
\newcommand{\samedef}{\ensuremath{\us {\text{def}} \Leftrightarrow}}
\newcommand{\qsamedef}{\quad\samedef\quad}
\newcommand{\lsamedef}{\ensuremath{\us {\text{def}} \Longleftrightarrow}}
\newcommand{\qlsamedef}{\quad\lsamedef\quad}


%%
% Mengendefinitionen
%%
\newcommand{\M}{\ensuremath{\mathcal{M}}}
\newcommand{\N}{\ensuremath{\mathbb{N}}}
\newcommand{\Z}{\ensuremath{\mathbb{Z}}}
\newcommand{\Q}{\ensuremath{\mathbb{Q}}}
\newcommand{\I}{\ensuremath{\mathbb{I}}}
\newcommand{\R}{\ensuremath{\mathbb{R}}}
\newcommand{\C}{\ensuremath{\mathbb{C}}}
\renewcommand{\S}{\ensuremath{\mathbb{S}}}
\newcommand{\Prim}{\ensuremath{\mathbb{P}}}
\newcommand{\0}{\ensuremath{\emptyset}}
\renewcommand{\Re}{\ensuremath{\operatorname{Re}}}
\renewcommand{\Im}{\ensuremath{\operatorname{Im}}}
\newcommand{\strich}{\ensuremath{\big\vert}}

%%
% Mengenoperation
%%
\renewcommand{\nin}{\not\in}


%%
% Logisches Und und Oder
%%
\newcommand{\sand}{\ensuremath{\; \land \;}}
\newcommand{\sor}{\ensuremath{\; \lor \;}}
\renewcommand{\*}{\ensuremath{\cdot}}
\newcommand{\oder}{\vee}
\newcommand{\n}{\neg}
\newcommand{\und}{\wedge}


%%
% Bei Integral, Summe,... die Grenzen über bzw.
% unter das Zeichen
%%
\renewcommand{\d}[1]{\ensuremath{\operatorname{d}\!{#1}}}
\newcommand{\Sum}{\sum\limits}
\newcommand{\Prod}{\prod\limits}
\newcommand{\Min}{\min\limits}
\newcommand{\Max}{\max\limits}
\newcommand{\Lim}{\lim\limits}
\newcommand{\Inf}{\inf\limits}
\newcommand{\Sup}{\sup\limits}
\newcommand{\Liminf}{\liminf\limits}
\newcommand{\Limsup}{\limsup\limits}
%\renewcommand{\Cup}{\bigcup\limits}
%\renewcommand{\Cap}{\bigcap\limits}
\newcommand{\Otimes}{\bigotimes\limits}
\newcommand{\Oplus}{\bigoplus\limits}


%%
% Arcus Cotangens, Logarithmus
%%
\let\oldsin\sin
\renewcommand{\sin}[1]{\oldsin{\left(#1\right)}}
\let\oldcos\cos
\renewcommand{\cos}[1]{\oldcos{\left(#1\right)}}
\let\oldtan\tan
\renewcommand{\tan}[1]{\oldtan{\left(#1\right)}}
\let\oldarcsin\arcsin
\renewcommand{\arcsin}[1]{\oldarcsin{\left(#1\right)}}
\let\oldarccos\arccos
\renewcommand{\arccos}[1]{\oldarccos{\left(#1\right)}}
\let\oldarctan\arctan
\renewcommand{\arctan}[1]{\oldarctan{\left(#1\right)}}
\let\oldln\ln
\renewcommand{\ln}[1]{\oldln{\left(#1\right)}}
\newcommand{\Log}[2]{\ensuremath{\operatorname{log}_{#1}\left(#2\right)}}
\newcommand{\e}{\operatorname{e}}
\renewcommand{\exp}[1]{\ensuremath{\operatorname{\e}^{#1}}}


%%
% Griechische Buchstaben
%%
\renewcommand{\epsilon}{\ensuremath{\varepsilon}}
\renewcommand{\phi}{\varphi}
\renewcommand{\l}{\lambda}
\renewcommand{\t}{\tau}


%%
% Bezeichnungen
%%
\newcommand{\Rgn}{\R_{>0}} % R grösser Null
\newcommand{\Rad}{\operatorname{Rad}}   % Radikal
\newcommand{\kgV}{\operatorname{kgV}}   % Kleinster gemeinsame Vielfache
\newcommand{\ggT}{\operatorname{ggT}}   % Grösster gemeinsame Teiler
\newcommand{\winkel}{\sphericalangle}          % Winkelsymbol
\DeclarePairedDelimiter\abs{\lvert}{\rvert} % besserer Betrag
\makeatletter
\let\oldabs\abs
\def\abs{\@ifstar{\oldabs}{\oldabs*}}
\makeatother


%%
% Vektorgeometrie
%%
\renewcommand{\v}[1]{\vv{#1}}       % besserer Pfeil über Buchstaben
\renewcommand{\vec}[2]{\begingroup\renewcommand{\arraystretch}{1}\left(\begin{array}{c} #1 \\ #2 \end{array}\right)\endgroup}   % two dimensional vector
\renewcommand{\Vec}[3]{\begingroup\renewcommand{\arraystretch}{1}\left(\begin{array}{c} #1 \\ #2 \\ #3 \end{array}\right)\endgroup}      % three dimensional vector


%%
% Text
%%
\newcommand*{\Ueq}[1]{\mathrel{\mathop{=}\limits_{\text{#1}}}}
\newcommand*{\Oeq}[1]{\mathrel{\stackrel{\text{#1}}{=}}}
\newcommand{\utext}[2]{\underbrace{#1}_{\substack{#2}}}
\newcommand{\otext}[2]{$\overbrace{\textrm{#1}}^{\textrm{#2}}$}
\newcommand{\lineortext}[1]{\hspace{0,1cm}\underline{\hspace{0.1cm}\hphantom{#1}\hspace{0.1cm}}\hspace{0,1cm}} % Lückentext erstellen
\newcommand{\gap}[1]{\rule{#1 cm}{1pt}}
%=========================================================
\newcommand{\theme}{Bruchterme}
\newcommand{\teacher}{Jonas Guler}
%=========================================================
\begin{document}
\fontsize{14}{14}\selectfont
\begin{center}
    \huge\bf\theme
\end{center}

\begin{questions}
%=========================================================
%-------------------TESTFRAGEN
%=========================================================
\question[4] %
Zerlege die Nenner der Brüche in seine Faktoren und gib den gemeinsamen Nenner an.
\begin{parcols}[2]
    \part $\frac{25}{a^2bc^3}$;\:$\frac{15}{ab^2c}$
    \part $\frac{12}{4p^2+4p+1}$;\:$\frac{6}{18p^3+9p^2}$
\end{parcols}

%---------------------------------------------------------
\question[4] %
Zerlege die Nenner der Brüche in seine Faktoren und gib den gemeinsamen Nenner an.
\begin{parcols}[2]
    \part $\frac{24}{x^2+7x+10}$;\:$\frac{48}{x^2+2x-15}$;\:$\frac{32}{x^2-x-6}$
    \part $\frac{4}{x^2-4y^2}$;\:$\frac{12}{x^2-4xy+4y^2}$;\:$\frac{3}{x^2-2xy}$
\end{parcols}

%---------------------------------------------------------
\question[4] %
Finde das kleinste gemeinsame Vielfache und den grössten gemeinsamen Teiler.
\begin{parcols}[2]
    \part $a^2bc^3$,\quad$ab^2c$
    \part $x^2-1$,\quad$x+1$
\end{parcols}

%---------------------------------------------------------
\question[4] %
Finde das kleinste gemeinsame Vielfache und den grössten gemeinsamen Teiler.
\begin{parcols}[2]
    \part $ax-a$,\quad$b-bx$,\quad$abx^2-ab$
    \part $x^2-2x+1$,\quad$x^2-1$,\quad$x-1$
\end{parcols}

%---------------------------------------------------------
\question[4]
Berechne das kleinste gemeinsame Vielfache und der grösste gemeinsame Teiler der folgenden beiden Termen.
\[4u^2-9v^2\quad\text{und}\quad 10u+15v\]

%---------------------------------------------------------
\question[4]
Berechne das kleinste gemeinsame Vielfache und der grösste gemeinsame Teiler der folgenden beiden Termen.
\[25a^2-16b^4\quad\text{und}\quad 15a^2+12ab^2\]


%---------------------- ADDITION -----------------------------------
\question[6]
Erweitere die Brüche auf einen gemeinsamen Nenner und vereinfache dann den Term so weit wie möglich.
\begin{parts}
    \part $\frac{z}{z-5}-\frac{5}{z+3}-\frac{40}{z^2-2z-15}$ %Sammlung 174a
    \part $\frac{m^2-8m}{(m-2)(2m-5)}-\frac{m}{5-2m}$ %Sammlung 178d
    \part $\frac{k+2}{6k-15}+\frac{8k+1}{8k-20}+\frac{k+11}{10-4k}$ %Sammlung 182a
\end{parts}

%---------------------------------------------------------
\question[6]
Erweitere die Brüche auf einen gemeinsamen Nenner und vereinfache dann den Term so weit wie möglich.
\begin{parts}
    \part $\frac{2x}{x-5}+\frac{3x}{x+5}$ %KSWillisau A2d
    \part $\frac{5w}{w-3}-\frac{3w}{w+2}-\frac{2w-1}{w^2-w-6}$ %KSWillisau A2f
    \part $\frac{a+3b}{2a-2b}+\frac{4a-2b}{a+b}-\frac{3a+4b}{2a+2b}$ %KSWillisau A2h
\end{parts}


%---------------------------------------------------------
\question[6]
Erweitere die Brüche auf einen gemeinsamen Nenner und vereinfache dann den Term so weit wie möglich.
\begin{parcols}[2]
    \part $\displaystyle\frac{3}{2x-1}+\frac{4}{4x-2}-\frac{5}{6x-3}$ %KSWillisau A2e
    \part $\displaystyle\frac{2x-8}{x^2-x-12}-\frac{3x-6}{x^2-5x+6}$ %KSWillisau A2k
\end{parcols}

%--------------------- MULTIPLICATION ------------------------------------
\question[6]
Multipliziere oder dividiere die Brüche und vereinfache dann den Term so weit wie möglich.
\begin{parts}
    \part $\frac{7r^2s}{12(r-2)}\cdot\frac{(2s-2r)^2}{21rs^2}$ %Sammlung 188c
    \part $\frac{x^2-xy}{x+y}:\frac{x-y}{3x+3y}$ %Sammlung 196b
    \part $\frac{c-5}{\frac{c^2-3c-10}{c^2-4}}$ %Sammlung 208b
\end{parts}

%---------------------------------------------------------
\question[6]
Multipliziere oder dividiere die Brüche und vereinfache dann den Term so weit wie möglich.
\begin{parts}
    \part $\frac{x^2+4x+4}{4z}\cdot\frac{8z}{x^2+5x+6}$ %KSWillisau A3e
    \part $\frac{48}{23b}:\left(-\frac{24b^5}{5}\right)$ %KSWillisau A4c
    \part $\frac{\frac{(x-3y)^2}{6xy}}{\frac{6y+2x}{xy}}$ %KSWillisau A4e
\end{parts}

%---------------------------------------------------------
\question[4]
Multipliziere die Brüche und vereinfache den Term so weit wie möglich.
\begin{parcols}[2]
    \part $\displaystyle\frac{4(x+2)}{5}\*\frac{25x}{x^2+4x+4}\*\frac{2x+4}{35}$ %KSWillisau A3c
    \part $\displaystyle\left(\frac{5p^2}{7y^2}-\frac{7q^2}{5y^2}\right)\*\left(\frac{35y^2}{5p+7q}\right)$ %KSWillisau A5b
\end{parcols}

%-------------------- MIXED -------------------------------------
\question[6]
Vereinfache die Terme so weit wie möglich.
\begin{parts}
    \part $\frac{25x^2-9}{(x+2)^2}\cdot\frac{x^2+5x+6}{y^3}:\frac{5x-3}{xy^3+2y^3}$ %Alg1 S108 A161b
    \part $\frac{\frac{x+h}{x+h-1}-\frac{x}{x-1}}{h}$ % DMK910 S34 A10d
    \part $\left(\frac{8x^2+4x+1}{4x^2-2x}-\frac{2x}{2x-1}\right)\cdot\frac{6x-3}{4x^2+2x}$ %DMK910 S34 A8b
\end{parts}

%---------------------------------------------------------
\question[6] %gemischt
Vereinfache die Terme so weit wie möglich.
\begin{parts}
    \part $\frac{4bd}{c}\cdot\left(\frac{3a}{4b}+\frac{15c}{6d}\right)$ %KSWillisau A5a
    \part $\frac{\frac{x+h}{x+h-1}-\frac{x}{x-1}}{h}$ % DMK910 S34 A10d
    \part $\left(\frac{8x^2+4x+1}{4x^2-2x}-\frac{2x}{2x-1}\right)\cdot\frac{6x-3}{4x^2+2x}$ %DMK910 S34 A8b
\end{parts}





\end{questions}
\end{document}